\section{Introduction}
\label{sec:intro}

Long-horizon LLM agents increasingly operate over weeks or months of
conversation, during which the \emph{facts} that describe a user change:
jobs end, preferences rotate, diagnoses update, plans get cancelled.
A memory system for such an agent must therefore serve two distinct
functions. It must \emph{find} relevant history (the retrieval problem
that dominates the Retrieval-Augmented Generation literature), and it
must also \emph{maintain truth} across time, so that the text it
resurfaces actually describes the user's current state rather than a
superseded one. The second capability is comparatively under-studied:
recent surveys of agent memory list ``memory automation'' and
``trustworthiness'' as open frontiers \citep{survey_agent_memory}, and
the benchmarks that stress these capabilities -- LongMemEval
\citep{longmemeval}, LoCoMo \citep{locomo}, and MemoryAgentBench
\citep{memoryagentbench} -- continue to show large gaps on questions
that require the system to know \emph{which fact is currently true}.

A-Mem \citep{amem} makes an important step away from flat embedding
stores by treating memory as a self-organising graph of structured
notes, each carrying keywords, context, and tags, with links that are
formed and rewritten as new notes arrive. This mechanism works well on
retrieval-centric question types but, because it operates at the level
of free-text notes, it has no explicit handle on \emph{when} a
statement was true, \emph{whether} a later statement replaced it, or
\emph{why} two statements might be incompatible. We show that this gap
is visible as a concrete failure mode on LongMemEval's
\emph{knowledge-update}, \emph{temporal-reasoning}, and
\emph{abstention} slices, and that closing it requires changing the
unit of memory, not the size of the reader.

\paragraph{Our approach.} We propose the \textbf{Temporal
Truth-Maintained Memory Graph} (TTMG). Instead of free-text notes, the
atomic memory unit is a \emph{claim}: a subject--predicate--object-like
assertion accompanied by a \emph{validity interval}, a
\emph{provenance} record (session id, turn id, speaker, session
timestamp), a \emph{polarity}, a \emph{volatility} class, and a
\emph{confidence}. When a new claim is ingested, the system retrieves
its semantic and lexical neighbours and uses a light LLM linker to
label each pair as \texttt{support}, \texttt{contradict},
\texttt{supersede}, or \texttt{unrelated}. \texttt{Supersede} edges
flip the older claim's \texttt{active} bit so that it is excluded from
retrieval for current-state questions but remains accessible for
historical queries. At read time, TTMG parses the question for its
temporal anchor and current/history intent, runs approximate
$k$-nearest-neighbour search on the \emph{active} subset, and then
extracts the \emph{maximum consistent subgraph} -- the largest set of
retrieved claims no two of which are linked by a hard
\texttt{contradict} edge. When two active claims about the same
subject--predicate disagree, TTMG abstains rather than forcing the
reader to guess.

\paragraph{Contributions.}
\begin{itemize}
  \item We introduce a \emph{claim-level} schema for long-horizon agent
  memory that couples semantic embedding indices with explicit
  validity intervals, provenance, polarity, and volatility
  (\S\ref{sec:method}).
  \item We introduce \emph{contradict} and \emph{supersede} edges, a
  lightweight LLM linker with a cheap subject--predicate gate, and an
  intent-gated hybrid reader that prefers claims for temporal
  questions and raw turns as a fallback on everything else. On a
  150-question stratified slice of LongMemEval-S, TTMG achieves
  paired set-inclusion wins over a strong flat hybrid-RAG baseline on
  two capability slices: \textbf{abstention} accuracy
  (\TTMGLMEAccABS\% vs \FlatLMEAccABS\%, +\DeltaLMEABS~pp, $n{=}9$,
  paired $b{=}4, c{=}0$; exact McNemar one-sided
  $p{=}$\McnemarABSonesided) and \textbf{single-session-user} accuracy
  (\TTMGLMEAccSSU\% vs \FlatLMEAccSSU\%, $n{=}21$, paired 2--0).
  Knowledge-update and temporal-reasoning are tied with the baseline
  at this scale (paired discordance symmetric, McNemar $p{=}1.0$),
  and we show that TTMG and Flat RAG answer the same slice accuracy
  through \emph{complementary} subsets of questions (\S\ref{sec:exp}).
  \item We propose \textbf{TTMG+R}, a type-aware router that composes
  TTMG with Flat RAG. Routing \texttt{single-session-user} traffic to
  TTMG and everything else to Flat RAG reaches \RouterLMEAcc\%
  overall, matching the per-slice best system and improving over
  Flat RAG by $+1.3$~pp while retaining TTMG's truth-maintenance
  behaviours (abstention, supersede) wherever they are live
  (\S\ref{ssec:lme}). An oracle upper bound of \OracleLMEAcc\%
  documents that claim-level and raw-text retrieval carry
  substantially complementary information.
  \item The mechanism is cheaper than the A-Mem note-graph baseline:
  TTMG uses \TTMGTokens{} tokens per question vs.\ \AmemTokens{}
  for A-Mem (\TokensOverhead\%) and runs in \TTMGLatency{}s vs.\
  \AmemLatency{}s, because the session-batched writer replaces
  A-Mem's per-turn LLM analysis (\S\ref{sec:efficiency}).
  \item An ablation of the LLM linker drops overall accuracy
  (-2.6~pp), abstention (-22~pp), knowledge-update (-9~pp),
  temporal-reasoning (-5~pp), and single-session-user (-5~pp),
  confirming that write-time truth maintenance -- not the claim
  schema alone -- supplies the gain (\S\ref{sec:ablation}).
\end{itemize}

The claim-level abstraction therefore turns memory from ``closest
text'' into ``currently true fact,'' delivers paired wins on the two
capability slices that truth maintenance should serve, composes with
Flat RAG to reach parity on the rest of the benchmark, and does so
at lower compute cost than a standard note-graph baseline.
